\documentclass[11pt]{article}
\usepackage[russian]{babel}
\usepackage{fontspec}

\setmainfont{Times New Roman}
\setsansfont{Arial}
\setmonofont{Courier New}

% 加载 TikZ 及其库
\usepackage{tikz}
\usetikzlibrary{graphs}

\begin{document}
	
	\section{Упражнение 8.2.1: Сложный граф в TikZ}
	
	\begin{center}
		\begin{tikzpicture}[
			% 定义节点样式
			
			%пределение стилей (.style): В начале окружения заданы шаблоны greenNode (зеленая заливка) и doubleNode (двойная граница), 
			%что позволяет не дублировать настройки для каждого узла.
			greenNode/.style={circle, draw, fill=green, minimum size=0.8cm},
			doubleNode/.style={circle, draw, double, double distance=1pt, minimum size=0.8cm}
			]
			% 1. 使用极坐标放置节点 (60度一个，半径3cm)
			%Полярные координаты: Узлы расставлены с помощью синтаксиса (угол:радиус).
			%Это позволило идеально распределить шесть точек по кругу с шагом $60^\circ$.
			\node[greenNode] (A) at (90:3cm) {A};
			\node[doubleNode] (B) at (150:3cm) {B};
			\node[greenNode] (C) at (210:3cm) {C};
			\node[doubleNode] (D) at (270:3cm) {D};
			\node[greenNode] (E) at (330:3cm) {E};
			\node[doubleNode] (F) at (30:3cm) {F};
			
			% 2. 绘制蓝色虚线并添加标签
			%Пунктирные линии и метки: Команды draw[dotted] соединяют узлы B, D и F. Метки (6, 2, 4) размещены с помощью параметра midway, который автоматически находит центр линии.
			\draw[blue, dotted, thick] (B) -- (F) node[midway, above] {6};
			\draw[blue, dotted, thick] (B) -- (D) node[midway, left] {2};
			\draw[blue, dotted, thick] (D) -- (F) node[midway, right] {4};
			
			% 3. 绘制红色曲线
			%Кривые Безье и слои: Линия от A к D реализована через .. controls .., что создает плавный S-образный изгиб. Математический символ $\sqrt{2}$ помещен в узел с белым фоном (fill=white), чтобы "перекрыть" саму линию для лучшей читаемости.
			%Изогнутые пути (bend): Для красных и синих дуг использован параметр to[bend left/right], который выгибает линию на заданный угол, не требуя ручного расчета координат.
			\draw[red] (A) to[bend right=40] (B);
			\draw[red] (A) to[bend left=40] (F);
			\draw[red] (F) to[bend left=40] (E);
			% A 到 D 的 S 形曲线
			\draw[red] (A) .. controls +(left:2cm) and +(right:2cm) .. (D) 
			node[pos=0.5, fill=white, inner sep=1pt] {$\sqrt{2}$};
			
			% 4. 绘制青蓝色实线曲线
			\draw[cyan, thick] (B) to[bend right=30] (C);
			\draw[cyan, thick] (C) to[bend right=30] (D);
			
		\end{tikzpicture}
	\end{center}
	
\end{document}